\subsection{Solutions des ateliers}
\begin{enumerate}
\item Le regroupement donne $M_n=n(6+(4n+2))/2=2n(n+2)$. Donc $M_8=160$. La fonction $M_n$ croît pour $n\geq0$ ; $M_{10}=240\leq260<M_{11}=286$ : on peut construire 10 couronnes complètes.
\item L'identité $1/((j+2)(j+3))=1/(j+2)-1/(j+3)$ donne $U_n=1/3-1/(n+3)$. Dans le produit, les facteurs intérieurs s'annulent et $V_n=(n+3)/3$. Ainsi $U_n=(1-1/V_n)/3$.
\item Avec k pièces B, choisir leurs positions puis leurs modèles donne $\binom5k3^{5-k}2^k$. Les six effectifs sont $243, 810, 1080, 720, 240, 32$ et leur total vaut 3125. L'autre comptage donne directement $5^5=3125$, car chaque emplacement possède cinq choix.
\end{enumerate}
